\documentclass[a4paper, 11pt]{article}

%\usepackage[
%backend=biber,
%style=numeric,
%]{biblatex}
\usepackage{physics}
%\addbibresource{reference.bib}

\begin{document}
\title{}
\author{}
\maketitle

\section*{Exercise 1}
Quantum phase estimation (QPE) can be used to measure the Hamming weight of a quantum state and prepare Dicke states.

\subsubsection*{a) Write down the Hamiltonian whose eigenvalues can distinguish the Hamming weight of an n qubit quantum state.}

The Hamiltonian whose eigenvalues differs by Hamming weight in $n$-qubit system can be constructed by the projector onto the space of fixed Hamming weight $i$. That is,
$$
\mathcal{H} = \sum_{i=0}^{2^n -1} e_i P_{i}
$$
We can choose $e_i = i$ where the eigenvalue becomes the Hamming weight $i$.
Note that the given hamiltonian can have ${}_{n}C_{i}$ for $e_i$.

\subsubsection*{b) For an $n$ qubit state, what is the minimum number of bits needed to measure the hamming weight of the state? Hint: The Hamming weight of the state is in the range $\{0,1,..,n\}$}
Let $k$ denote the minimum number of bits required to measure hamming weight of the state. Then
$$
k = \left[\log_{2}(n+1)\right]
$$
where $\left[ a \right]$ denotes the smallest integer larger than $a$.

\section*{Exercise 2}
Using qiskit, create a function that takes in a quantum state in vector format as input and returns a quantum circuit that measures the Hamming weight of the state. Hint: see ref \cite{rethinasamy2024logarithmicdepthquantumcircuitshamming}. You can use qiskit's initialize module to initialize the states.

\section*{Exercise 3}
Use the function created in part 2 above to measure the Hamming weight of the following states:

\subsubsection*{a) $\ket{11}$}
\subsubsection*{b) $\ket{01}$}
\subsubsection*{c) $\frac{1}{\sqrt{2}} \left(\ket{00} + \ket{11} \right)$}
\subsubsection*{d) $\frac{1}{\sqrt{3}} \left( \ket{011} + \ket{101} + \ket{110} \right)$}
\subsubsection*{e) Explain the result in part c) above.}

\section*{Exercise 4}
Explain how the function in part 2 above can be used to prepare a Dicke state.

\section*{Exercise 5}
Suppose we want to prepare the Dicke state with $n = 100$ and $k = 50$.

\subsubsection*{What is the probability that the state will be measured?}

\subsubsection*{What is the minimum number of measurements required?}

\subsubsection*{How many gates are required?}

%\printbibliography


\bibliography{reference}
\bibliographystyle{plain}

\end{document}